\documentclass[10pt,a4paper]{amsart}
\usepackage{amsmath,amssymb,amsthm,graphicx}
\usepackage{verbatim}
\newenvironment{reference}{\comment}{\endcomment}
%\newenvironment{reference}{}{}
\newenvironment{slogan}{\comment}{\endcomment}
\newenvironment{history}{\comment}{\endcomment}
% For commutative diagrams we use Xy-pic
\usepackage[all]{xy}

% We use 2cell for 2-commutative diagrams.
\xyoption{2cell}
% Theorem environments.
%
\theoremstyle{plain}
\newtheorem{theorem}[subsection]{Theorem}
\newtheorem{proposition}[subsection]{Proposition}
\newtheorem{lemma}[subsection]{Lemma}

\theoremstyle{definition}
\newtheorem{definition}[subsection]{Definition}
\newtheorem{example}[subsection]{Example}
\newtheorem{exercise}[subsection]{Exercise}
\newtheorem{situation}[subsection]{Situation}

\theoremstyle{remark}
\newtheorem{remark}[subsection]{Remark}
\newtheorem{remarks}[subsection]{Remarks}

\numberwithin{equation}{subsection}

% Macros
%
\def\lim{\mathop{\mathrm{lim}}\nolimits}
\def\colim{\mathop{\mathrm{colim}}\nolimits}
\def\Spec{\mathop{\mathrm{Spec}}}
\def\Hom{\mathop{\mathrm{Hom}}\nolimits}
\def\Ext{\mathop{\mathrm{Ext}}\nolimits}
\def\SheafHom{\mathop{\mathcal{H}\!\mathit{om}}\nolimits}
\def\SheafExt{\mathop{\mathcal{E}\!\mathit{xt}}\nolimits}
\def\Sch{\mathit{Sch}}
\def\Mor{\mathop{\mathrm{Mor}}\nolimits}
\def\Ob{\mathop{\mathrm{Ob}}\nolimits}
\def\Sh{\mathop{\mathit{Sh}}\nolimits}
\def\NL{\mathop{N\!L}\nolimits}
\def\CH{\mathop{\mathrm{CH}}\nolimits}
\def\proetale{{pro\text{-}\acute{e}tale}}
\def\etale{{\acute{e}tale}}
\def\QCoh{\mathit{QCoh}}
\def\Ker{\mathop{\mathrm{Ker}}}
\def\Im{\mathop{\mathrm{Im}}}
\def\Coker{\mathop{\mathrm{Coker}}}
\def\Coim{\mathop{\mathrm{Coim}}}

% Boxtimes
%
\DeclareMathSymbol{\boxtimes}{\mathbin}{AMSa}{"02}

%
% Macros for moduli stacks/spaces
%
\def\QCohstack{\mathcal{QC}\!\mathit{oh}}
\def\Cohstack{\mathcal{C}\!\mathit{oh}}
\def\Spacesstack{\mathcal{S}\!\mathit{paces}}
\def\Quotfunctor{\mathrm{Quot}}
\def\Hilbfunctor{\mathrm{Hilb}}
\def\Curvesstack{\mathcal{C}\!\mathit{urves}}
\def\Polarizedstack{\mathcal{P}\!\mathit{olarized}}
\def\Complexesstack{\mathcal{C}\!\mathit{omplexes}}
% \Pic is the operator that assigns to X its picard group, usage \Pic(X)
% \Picardstack_{X/B} denotes the Picard stack of X over B
% \Picardfunctor_{X/B} denotes the Picard functor of X over B
\def\Pic{\mathop{\mathrm{Pic}}\nolimits}
\def\Picardstack{\mathcal{P}\!\mathit{ic}}
\def\Picardfunctor{\mathrm{Pic}}
\def\Deformationcategory{\mathcal{D}\!\mathit{ef}}
\usepackage{cleveref}
\crefname{originalreference}{original reference}{original references}
\Crefname{originalreference}{Original reference}{Original references}
\makeatletter
\newlabel{derived-lemma-grothendieck-spectral-sequence}{{derived-lemma-grothendieck-spectral-sequence}{source}{}{sourcepointer-1}{}}
\newlabel{derived-lemma-grothendieck-spectral-sequence@cref}{{[originalreference][0][]derived-lemma-grothendieck-spectral-sequence}{[1][0][]source}{}{}{}}
\makeatother
\title{A Grothendieck spectral sequence for two specialization-stable supports and the composition of the associated derived functors}
\author{The Stacks Project Authors}
\date{stacks-project:local-cohomology.tex:local-snapshot-20261012; GNU FDL 1.2 or later}
\begin{document}
\maketitle
\noindent\textit{Editorial note.} All delivered mathematical text is the author's own, copied from the pinned Stacks Project chapter source sources/stacks-local-cohomology/pinned-original-source.tex; no statement or proof step was rewritten, reordered or removed.\par\medskip
\noindent\textit{Editorial note.} Preamble provenance: the Stacks Project keeps one single shared preamble for every chapter; this card takes it directly from the staged sources/stacks-local-cohomology/preamble.tex as three verbatim spans (lines 8-12, 14-18 and 154-217), so no derived file was written and nothing was added or redefined. The wrapper class and the base packages are supplied by the assembler.\par\medskip
\noindent\textit{Editorial note.} Bibliography provenance: neither the counted statement nor its proof contains a citation (this lane rescans the whole span before writing the card), so the card carries no bibliography and no entry was reconstructed.\par\medskip
\noindent\textit{Editorial note.} Selection provenance: this card was opened under the file-based control contract (CONTROL.md), with the mechanical pre-card gate run on its statement and proof interval before the card was written.\par\medskip
\noindent\textit{Allowed standard prerequisites (editorial).}\begin{itemize}
\item The functor R\textbackslash{}Gamma\_T of sections with support in a specialization-stable subset T and its right adjointness to the inclusion of the category of modules supported in T. The author's section on more general supports, in which the counted statement occurs. Identifies the composite of the two support functors with the functor for the intersection of the supports.
\item The Grothendieck spectral sequence of a composition of derived functors, in the form recorded in the chapter on derived categories. Invoked twice in the counted proof (statement line 1138-1139 and proof line 1158-1159) as a cross-chapter pointer. Produces the spectral sequence from the identification R\textbackslash{}Gamma\_\{T \textbackslash{}cap T'\} = R\textbackslash{}Gamma\_T \textbackslash{}circ R\textbackslash{}Gamma\_\{T'\}.
\end{itemize}\medskip
\ifdefined\hypertarget\hypertarget{sourcepointer-1}{}\fi
\noindent\textit{Source reference.} Original reference derived-lemma-grothendieck-spectral-sequence: sources/stacks-local-cohomology/pinned-original-source.tex, source lines 1139--1139. Navigation pointer only: the counted text refers to the Stacks item derived-lemma-grothendieck-spectral-sequence, which belongs to another chapter of the project and is not part of this card. The pointer carries the target's own label as its number because this chapter's staged snapshot does not contain the other chapter's tag table; its content is not reproduced here.\par\medskip
\begin{lemma}
\label{lemma-local-cohomology-ss}
Let $A$ be a ring and let $T, T' \subset \Spec(A)$ subsets
stable under specialization. For $K \in D^+(A)$
there is a spectral sequence
$$
E_2^{p, q} = H^p_T(H^p_{T'}(K)) \Rightarrow H^{p + q}_{T \cap T'}(K)
$$
as in Derived Categories, Lemma
\ref{derived-lemma-grothendieck-spectral-sequence}.
\end{lemma}

\begin{proof}
Let $E$ be an object of $D_{T \cap T'}(A)$. Then we have
$$
\Hom(E, R\Gamma_T(R\Gamma_{T'}(K))) =
\Hom(E, R\Gamma_{T'}(K)) =
\Hom(E, K)
$$
The first equality by the adjointness property of $R\Gamma_T$
and the second by the adjointness property of $R\Gamma_{T'}$.
On the other hand, if $J^\bullet$ is a bounded below complex
of injectives representing $K$, then $H^0_{T'}(J^\bullet)$
is a complex of injective $A$-modules representing $R\Gamma_{T'}(K)$
and hence $H^0_T(H^0_{T'}(J^\bullet))$ is a complex representing
$R\Gamma_T(R\Gamma_{T'}(K))$. Thus $R\Gamma_T(R\Gamma_{T'}(K))$
is an object of $D^+_{T \cap T'}(A)$. Combining these two
facts we find that $R\Gamma_{T \cap T'} = R\Gamma_T \circ R\Gamma_{T'}$.
This produces the spectral sequence by the lemma referenced
in the statement.
\end{proof}

\end{document}
